\section{Introduction}

Financial asset markets are dynamic and complex systems where any sudden change in the terms of trade can lead to significant price fluctuations, reduced liquidity and even systemic crises.
Such phenomena, referred to as market shocks, attract the attention of researchers and practitioners, as their consequences can have far-reaching effects on the economy as a whole.
In the current conditions of globalization and digitalization of financial markets, the issue of modeling market shocks becomes particularly relevant.

The purpose of this study is to develop and apply agent-based simulation to analyze the behavior of the financial system when market shocks occur.
The research uses a simulator implemented in the AgentBasedModel module to reproduce the dynamics of trading on exchanges involving traders, market makers and other agents.
The objectives of the research include:

\begin{enumerate}
    \item Studying the baseline behavior of the model without shocks to establish control dynamics.
    \item Introducing different scenarios of market shocks (e.g., a sudden change in liquidity or a sharp decline in demand) and analyzing their impact on key system indicators.
    \item Assessing the sensitivity of the model to changes in parameters such as the number of trading participants, intensity of shock events and the characteristics of market makers.
    \item Comparison of the obtained results with theoretical expectations and empirical data presented in the literature.
\end{enumerate}

The relevance of the study is determined by the growing need to develop effective forecasting and risk management tools in conditions of increased market volatility.
The results of the work can be useful for the formation of strategies to stabilize financial markets and the development of early warning methods for systemic risks.

The structure of the paper is as follows: the section “Literature Review” includes the analysis of the main approaches and comparative analysis of existing methods of modeling market shocks.
The main chapters contain a detailed description of the modeling methodology (Chapter 1), description of the experimental part and analysis of the obtained results (Chapter 2).
The conclusion summarizes the results of the study, discusses its achievements and prospects for further research.